\documentclass[11pt]{article}
\usepackage[margin=25mm]{geometry}
\usepackage{fontspec}\setmainfont{Latin Modern Roman}
\usepackage{amsmath,amssymb,mathtools}
\usepackage{unicode-math}\setmathfont{Latin Modern Math}
\usepackage{xurl,hyperref}
\hypersetup{colorlinks=true,urlcolor=blue}
\setcounter{secnumdepth}{0}\setlength{\parindent}{0pt}\setlength{\parskip}{5pt}\setlength{\emergencystretch}{3em}
\newcommand{\sref}[2]{\hyperlink{source:#1:#2}{[S#2]}}
\newcommand{\cataloguescope}[1]{\paragraph{Recorded target.}#1}
\begin{document}
\section{Kourovka Notebook Problem 21.72}\label{statement}

\subsection{Definitions and mathematical statement}

\paragraph{Shared notation.}$\mathbb N=\{1,2,\ldots\}$, and $\mathbb Z,\mathbb Q,\mathbb R,\mathbb C$ denote the integers, rationals, reals and complex numbers. Any other conventions are specified locally.


Use the source's definition of a Tarski monster: a finitely generated noncyclic group whose proper nontrivial subgroups are all isomorphic to each other. No restriction that these subgroups be cyclic is silently added. A discrete group is amenable if there is a positive linear functional $m:\ell^\infty(G)\to\mathbb R$ with $m(1)=1$ invariant under all left translations. A group is torsion-free when no nonidentity element has finite order, and has prime exponent $p$ when $g^p=1$ for every element and $G$ is nontrivial.
\paragraph{Statement recorded in the supplied source.}
(a) Does there exist an amenable torsion-free Tarski monster in the specified sense?

(b) Does there exist an amenable Tarski monster of prime exponent? The prime is not fixed in advance. \sref{statement}{2}
\subsection{Short English statement}
Do amenable Tarski monsters exist, either torsion-free or of prime exponent?
\subsection{Sources}
\begin{description}
\item[\hypertarget{source:statement:1}{S1}] ULAM.AI and the UnsolvedMath Contributors, \emph{UnsolvedMath: A Curated Collection of Open Mathematics Problems}, catalogue entry. CC BY 4.0. \url{https://huggingface.co/datasets/ulamai/UnsolvedMath}.
\item[\hypertarget{source:statement:2}{S2}] D. Kielak, Problem 21.72, in E. I. Khukhro and V. D. Mazurov (eds.), \emph{The Kourovka Notebook: Unsolved Problems in Group Theory}, 21st edition, arXiv:1401.0300v48 (6 October 2026). \url{https://arxiv.org/abs/1401.0300v48}.
\end{description}
\label{end:statement}
\end{document}
